\section{Method}
\label{sec:method}

Figure~\ref{fig:arch} shows the loop.
A trust head turns the current graph into one scalar $\text{trust}_v\in[0,1]$ per RSU, and nothing else; a fixed rule uses it to pick a target; the outcome of the task updates the behavioural features that the next step's graph carries.
% src: A1, L1

\begin{figure*}[t]
\centering
\includegraphics[width=0.92\textwidth]{\figdir/fig1_architecture.pdf}
\caption{System overview: dynamic graph, trust head, fixed offloading rule, task outcome fed back as features, and a read-only explanation. The sealed ground truth is used only to evaluate. Ablation variants A--E use the same rule. (Generated by \texttt{scripts/make\_figures.py}.)}
\label{fig:arch}
\end{figure*}

\subsection{Trust heads C and D}
Both heads are trained on the identical self-supervised objective and read the same RSU feature rows; the difference is whether the head also aggregates over neighbours (with the edge features) by message passing.
% src: PROJECT_SPEC 3 (C and D trained on identical objective and features); A2
\emph{C} is a node-local MLP over the RSU's own six feature values and never sees a neighbour.
\emph{D} is a 2-layer GINEConv over the graph, with message passing along all four edge features; C is matched to D in depth and width.
% src: A4 (C node-local MLP; D 2-layer GINEConv); D44 (C matches D's depth and width); A9 ("capacity-matched node-local head")
\TODO{cite GINEConv (not in the verified reference set of RELATED\_WORK.md)}
Both emit trust through a sigmoid.
The target is next-task success: whether the RSU completes the next task assigned to it within its deadline.
The loss is binary cross-entropy, and the injected behaviour class never enters a loss or a feature vector.
% src: A2 (L3, L4)
Each head is trained once per seed on mixed conditions ($\rho\in\{0,0.25,0.5,0.75,1\}$, degraded fraction $\in\{0.1,0.2,0.3\}$, injection draws $\{0,1\}$; validation and held-out draws $40$ and $90$, all draw indices below $100$), with Adam, learning rate $10^{-3}$, batch size 32, at most 200 epochs, patience 8, and model selection on held-out validation loss only.
% src: D43 section 2 and D68 section 2 (training draws, TrainSettings = D45: lr 1e-3, Adam, batch 32, cap 200, patience 8)
The training dispatch is the no-trust baseline (open loop), and the heads are then frozen and used unchanged across every evaluation condition.
% src: A2 (trained once per seed, frozen, L12); A19 (trained on Baseline-A open-loop dispatch)

\subsection{The fixed weighted rule}
\label{sec:rule}
For a task, the candidates are the in-range, dwell-feasible RSUs, and the target is
\begin{equation}
v^\star=\arg\max_v\;\alpha\,\text{trust}_v-\beta\,\text{latency}_v-\gamma\,\text{load}_v .
\end{equation}
% src: A1
Selection is not learned.
The weights $(\alpha,\beta,\gamma)$ were chosen once per variant on a separate tuning seed by a full-factorial grid and applied unchanged to every evaluation seed; the surface is flat for the trust-bearing variants, so the specific weights carry no claim.
% src: A18 (D50, D52, F20); A4 ("D52 weights, tuned once on seed 100 and applied unchanged")
Trust enters as a \emph{weighted term}, not as an admission gate: a moderately trusted but much closer node can still win.
This is a design argument; Sec.~\ref{sec:gate} reports what the gate comparison measured.
% src: A15 (D62 framing: design argument, not a measured gain); X3

\subsection{Closed-loop dispatch}
In the closed loop every task is dispatched by the variant under test, using the trust computed from the graph at that step; the observed outcome updates the RSU's \texttt{success\_ewma} and \texttt{latency\_dev}, so the dispatch changes the evidence that later steps see.
% src: A1, A3 (L8), Figure 1; D51
Training used open-loop baseline dispatch, so the closed loop is a distribution shift for the trust heads; we state it as a limitation (Sec.~\ref{sec:limits}).

\subsection{Variants and baselines}
\label{sec:variants}
Every variant uses the rule above; the variants differ only in the trust term.
% src: A1, A4; PROJECT_SPEC 3, L1/L2
\begin{itemize}
\item \textbf{A}: no trust, the same rule with $\alpha=0$. There is no learned selector anywhere in the project.
\item \textbf{B}: certificate-only trust. In this scenario \texttt{cert\_valid} is constant, so B makes the same decisions as A (asserted on all 2000 paired draws).
B is therefore \emph{not} a certificate system that was beaten; the comparison throughout is against no behavioural signal.
% src: A4 (F26 check), X4
\item \textbf{B-oracle}: trust read from the sealed fault state with zero delay. It is an \emph{upper bound}, not a method.
\item \textbf{C}: node-local MLP trust. \textbf{D}: 2-layer GINEConv trust over the graph.
\item \textbf{E}: D plus the explanation of Sec.~\ref{sec:explanation}. Its decisions equal D's by construction (asserted), so E is not a separate performance result.
% src: A4, X12
\item \textbf{Gate-C} (extra baseline): C's frozen head with candidates below $\tau=0.9$ removed first, then the same rule. If no candidate survives, the decision falls back to the ungated candidate set.
% src: A4, A15, D59 (fallback to the ungated set)
\item \textbf{H} (extra baseline): trust is the RSU's own \texttt{success\_ewma}; no model, no learning, no neighbour information, the same rule, no gate. Its weights $(0.5,0.25,0.3)$ were tuned once on the tuning seed by the same procedure as the other variants and applied unchanged.
% src: A4, A23 (H weights 0.5, 0.25, 0.3; F32; D65)
\item \textbf{D2} (exploratory, pre-registered): a contained graph head, C's node-local branch plus a mean over same-segment RSU neighbours only (no vehicle nodes), with neighbour weights initialised to zero so that it starts as C; 1570 parameters against C's 1313.
% src: A28; D68 section 1
\end{itemize}
The staggered-onset experiment (Sec.~\ref{sec:graph}) trains fresh heads of C, D and D2 on the staggered regime and gives each model its own weights tuned on the tuning seed.
% src: A25 (heads trained fresh on the staggered regime), A28
